\documentclass[11pt]{amsart}
\usepackage[margin=1in]{geometry}
\usepackage{amsmath,amssymb,amsthm}
\usepackage{enumitem}
\usepackage{listings}
\usepackage{xcolor}
\usepackage[colorlinks=true,linkcolor=blue!50!black,citecolor=blue!50!black,urlcolor=blue!50!black]{hyperref}

\newtheorem{theorem}{Theorem}[section]
\newtheorem{lemma}[theorem]{Lemma}
\newtheorem{proposition}[theorem]{Proposition}
\newtheorem{corollary}[theorem]{Corollary}
\newtheorem*{maintheorem}{Main Theorem}
\newtheorem*{conjecture*}{Conjecture (Z.-W. Sun, 2010)}
\theoremstyle{definition}
\newtheorem{definition}[theorem]{Definition}
\theoremstyle{remark}
\newtheorem{remark}[theorem]{Remark}

\newcommand{\Z}{\mathbb{Z}}
\newcommand{\Q}{\mathbb{Q}}
\newcommand{\F}{\mathbb{F}}
\newcommand{\Zp}{\mathbb{Z}_p}
\newcommand{\CT}{\operatorname{CT}}
\newcommand{\cF}{\mathcal{F}}
\newcommand{\cL}{\mathfrak{L}}
\newcommand{\bin}[2]{\binom{#1}{#2}}

\lstset{basicstyle=\ttfamily\scriptsize,breaklines=true,columns=fullflexible,frame=single,
  keepspaces=true,showstringspaces=false,upquote=true}

\title[Sun's conjecture for $\sum(20k+17)W_k(1,-4)$]{Sun's divisibility conjecture for
$\sum_{k<n}(20k+17)\,W_k(1,-4)$ is true}
\subjclass[2020]{11A07, 11B65, 11F33, 14F30, 33C20}
\keywords{Sun's conjectures, Dwork crystals, unit-root crystal, Ramanujan series for $1/\pi$,
complex multiplication, supercongruences}
\author{QEEDDD Verifiable-Mathematical-Reasoning-Corpus project}
\date{November 8, 2026}

\begin{document}

\begin{abstract}
Let $a_n=\sum_{k}(-4)^k\binom nk^2\binom{n-k}k^2=W_n(1,-4)$ and
$S(n)=\sum_{k=0}^{n-1}(20k+17)a_k$. In 2010 Z.-W.~Sun conjectured that $n\mid S(n)$ for every
$n\ge1$ (OEIS A179524). We prove this conjecture. Earlier work established the $2$-part, the
$5$-part and all prime powers $p^e$ with $p\equiv11,13,17,19\pmod{20}$. The split primes
$p\equiv1,3,7,9\pmod{20}$ remained open. At these primes the complex-multiplication elliptic curve behind the
problem has ordinary reduction, and none of the earlier vanishing-moment mechanisms applies.
We treat them with the unit-root crystals of Beukers and Vlasenko. First, an exact identity
reduces $S(N)\bmod N$ to the constant term of $\Lambda^N\rho\Phi(uv)$. This reduces the $p$-part of
the conjecture to a single $p$-adic number: the value at $b=-4$ of the Gauss--Manin connection
$\eta_b$ on a rank-one unit-root crystal. Our main new result identifies this connection:
$\eta_b(b)=-\tfrac12-\eta_{\mathcal R}(1/(256b))$, where $\eta_{\mathcal R}$ is the
unit-root connection of Ramanujan's family
${}_3F_2(\tfrac14,\tfrac12,\tfrac34;1,1;256T)$. The proof expands a second
model of the crystal at a corner of its Newton polytope; the resulting Hasse--Witt coefficients
converge $p$-adically to Ramanujan's period $c^{-1}\,{}_3F_2(\tfrac14,\tfrac12,\tfrac34;1,1;1/c^{2})$.
Sun's theorem on the partial sums of Ramanujan's series
$\sum(20k+3)\binom{2k}k^2\binom{4k}{2k}(-2^{10})^{-k}=8/\pi$ then gives
$20\eta_b(-4)+7=0$, which yields $p^{v_p(N)}\mid S(N)$ for every $N$. All steps are unconditional;
an exact-arithmetic verification script is included.
\end{abstract}

\maketitle

\tableofcontents

\section{Introduction}

For $n\ge0$ put
\[
 M(n,k)=\binom nk\binom{n-k}k=\binom n{2k}\binom{2k}k,\qquad
 a_n=\sum_{0\le2k\le n}(-4)^kM(n,k)^2,
\]
and
\[
 S(n)=\sum_{k=0}^{n-1}(20k+17)a_k .
\]
In Sun's notation $a_n=W_n(1,-4)$, where $W_n(b,c)=\sum_k\binom n{2k}^2\binom{2k}k^2b^{n-2k}c^k$.
The first values are $a_0,a_1,a_2,a_3=1,1,-15,-143$ and $S(1),\dots,S(4)=17,54,-801,-11812$.

\begin{conjecture*}[{\cite[Conjecture 5.8(ii)]{Sun2010}, \cite{OEIS}}]
For every positive integer $n$ we have $n\mid S(n)$.
\end{conjecture*}

The conjecture was stated in the eleventh version of \cite{Sun2010} (dated 1~July~2010 in
OEIS A179524 \cite{OEIS}). It is recorded as an open statement in the Lean database
\texttt{formal-conjectures} \cite{FC} and in the \emph{OEIS Open} benchmark \cite{OO}.

\begin{maintheorem}
For every integer $n\ge1$, $n\mid S(n)$. Moreover $v_2(S(n))=v_2(n)$.
\end{maintheorem}

\subsection{What was known}
The following facts are reproved below for completeness.
(i) $\mathrm{rad}(n)\mid S(n)$ for every $n$, by a Lucas congruence and a reflection that turns
$S(p)\bmod p$ into a truncation of Ramanujan's series \cite{EarlierNote,K26}.
(ii) $S(n)/n$ is an odd $2$-adic integer \cite{K26}.
(iii) $p^e\mid S(p^em)$ for all $e,m\ge1$ when $p=5$ or $p\equiv11,13,17,19\pmod{20}$
\cite{K26}. This uses an exact ``digit operator'' and the vanishing of three
moments, which comes from the supersingular reduction of the CM curve $j=j(\sqrt{-5})$.
At the \emph{split primes}, $p\equiv1,3,7,9\pmod{20}$, the CM curve is ordinary and the moments
are units. Before this paper only finitely many prime-power levels were known there,
by computation \cite{EarlierNote,K26}.

\subsection{The new result}
\begin{theorem}\label{thm:B}
Let $p$ be a prime with $p\equiv1,3,7,9\pmod{20}$. Then $p^{v_p(N)}\mid S(N)$ for every $N\ge1$.
\end{theorem}

The Main Theorem follows from Theorem~\ref{thm:B} together with Theorems~\ref{thm:two},
\ref{thm:inert} and~\ref{thm:five} (Section~\ref{sec:main}).

\subsection{Strategy for the split primes}
\begin{enumerate}[leftmargin=*]
\item \emph{An exact reduction} (Section~\ref{sec:elementary}). A telescoping identity gives
$S(N)=10N\cdot(\text{integer})+H_2(N)$, and $H_2(N)=\CT(\Lambda^N\rho\Phi(uv))$ with
$\Lambda=(1+u)(1+v)/(uv)$, $\rho=uv/(1+u+v)$ and
$\Phi(x)=\sum_l(20l+7)(-4)^l\binom{2l}l^2x^{2l}$. So it suffices to show that the power series
$\rho\Phi(uv)$ is a $p$-adic \emph{formal divergence} $\theta_uJ_1+\theta_vJ_2$.
\item \emph{A rank-one unit-root crystal} (Section~\ref{sec:crystal}). The series
$\rho\Phi(uv)$ is a constant term of the rational function
$\Omega=(20\theta_b+7)\bigl(uv/f_b\bigr)$ at $b=-4$, where
$f_b=(1+u+v)(1-uv(z+b/z)(w+1/w))$. At an ordinary prime, the Beukers--Vlasenko theory of Dwork
crystals \cite{BV1} gives $\Omega\equiv(20\eta_b+7)\,uv/f_b$ modulo formal derivatives,
with a unit-root connection $\eta_b$. So everything reduces to the single identity
$20\eta_b(-4)+7=0$.
\item \emph{Identification with Ramanujan's family} (Section~\ref{sec:identification}). We prove
$\eta_b(b)=-\frac12-\eta_{\mathcal R}(1/(256b))$, where $\eta_{\mathcal R}=\theta_T\log
\cF(T)$ and $\cF(T)=\sum_k\frac{(4k)!}{k!^4}T^k$. To prove it we use a second model
$f_c=(1-cxg(z,w))(1-xL(y))$ of the same crystal ($b=c^2$) and expand it at a
\emph{corner} vertex of its Newton polytope. There the Hasse--Witt coefficients are explicit
binomial sums that converge $p$-adically, as the level grows, to Ramanujan's period
$c^{-1}\cF(1/(256c^2))$.
\item \emph{Sun's theorem} (Section~\ref{sec:ramanujan}). Sun proved via a WZ pair that
$4(2n+1)\binom{2n}n$ divides the partial sums of Ramanujan's $(20k+3)$-series
\cite[Theorem~1.2]{SunEJC}. Together with Dwork's congruences \cite{MV,BV2} this gives
$\eta_{\mathcal R}(-1/1024)=-3/20$. Hence $20\eta_b(-4)+7=20(-\frac12+\frac3{20})+7=0$.
\end{enumerate}

\begin{remark}[An upside-down Ramanujan formula]\label{rem:upside}
The ``period'' of the crystal in step~(2) is
$\int_0^{1/4}{}_2F_1(\frac12,\frac12;1;16bx^2)\,(1-4x)^{-1/2}\,dx=\frac12\,{}_3F_2(\frac12,\frac12,\frac12;\frac34,\frac54;b)$,
a series that diverges at $b=-4$. The point of maximal unipotent monodromy of its differential
equation is $b=\infty$, where the holomorphic solution is
$b^{-1/2}{}_3F_2(\frac14,\frac12,\frac34;1,1;1/b)$. Formally, Sun's sum is the $p$-adic shadow of
the ``upside-down'' series $\sum_l(20l+7)(-4)^l\frac{(1/2)_l^3}{(1)_l(3/4)_l(5/4)_l}$. Up to sign,
this is Ramanujan's series $\sum_k(20k+3)\frac{(1/2)_k(1/4)_k(3/4)_k}{k!^3}(-\frac14)^k=\frac8\pi$
continued to the half-integers $k=-l-\frac12$. Theorem~\ref{thm:identification} makes this precise
$p$-adically, at every ordinary prime.
\end{remark}

\subsection*{Conventions}
$\theta_y=y\,\partial/\partial y$ for any variable $y$. $\CT$ denotes the constant term.
$v_p$ is the $p$-adic valuation, with $v_p(0)=\infty$. $(a)_k$ is the rising factorial. For an odd
prime $p$ we write $h=(p-1)/2$ and $d=\lfloor p/4\rfloor$. We put
$\gamma_l=(-4)^l\binom{2l}l^2$, $A_k=\frac{(4k)!}{k!^4}=\binom{2k}k^2\binom{4k}{2k}$, and
\[
 G_b(x)=\sum_{l\ge0}b^l\binom{2l}l^2x^{2l},\qquad
 E(x)=\sum_{i\ge0}\binom{2i}ix^i=(1-4x)^{-1/2},\qquad
 \Phi(x)=\sum_{l\ge0}(20l+7)\gamma_lx^{2l}.
\]
Thus $G_{-4}(x)=\sum_l\gamma_lx^{2l}$ and $\Phi=\bigl((20\theta_b+7)G_b\bigr)|_{b=-4}$.
For $N\ge1$ we put
\begin{equation}\label{eq:dN}
 d_N(b)=[x^{N-1}]\,G_b(x)E(x)=\sum_{0\le2l\le N-1}b^l\binom{2l}l^2\binom{2i}i,\qquad i=N-1-2l.
\end{equation}

\section{Elementary reductions}\label{sec:elementary}

\begin{lemma}\label{lem:telescope}
For integers $m,N\ge0$,
$\displaystyle\sum_{n=0}^{N-1}(2n-m+1)\binom nm^2=N\binom N{m+1}\binom{N-1}m$
(the right side is read as $0$ for $N=0$).
\end{lemma}
\begin{proof}
Put $g(n)=(n-m)^2\binom nm^2/(m+1)$. For $n\ge1$ this equals $n\binom n{m+1}\binom{n-1}m$, and
$g(0)=0$. Since $\binom{n+1}m=\frac{n+1}{n+1-m}\binom nm$ for $n\ge m$,
\[
 g(n+1)-g(n)=\binom nm^2\frac{(n+1)^2-(n-m)^2}{m+1}=(2n-m+1)\binom nm^2.
\]
This holds for all $n\ge0$; for $n<m$ both sides vanish. Now sum over $0\le n<N$.
\end{proof}

\begin{proposition}\label{prop:split}
Let
\[
 H_3(N)=\sum_{n<N}\sum_k(2n-2k+1)\gamma_k\binom n{2k}^2,\qquad
 H_2(N)=\sum_{n<N}\sum_k(20k+7)\gamma_k\binom n{2k}^2 .
\]
Then $S(N)=10H_3(N)+H_2(N)$ and
$H_3(N)=N\sum_k\gamma_k\binom N{2k+1}\binom{N-1}{2k}$. In particular $S(N)\equiv H_2(N)\pmod N$.
\end{proposition}
\begin{proof}
We have $a_n=\sum_k\gamma_k\binom n{2k}^2$ and $10(2n-2k+1)+(20k+7)=20n+17$. Apply
Lemma~\ref{lem:telescope} with $m=2k$.
\end{proof}

\begin{proposition}\label{prop:CT}
Let $\Lambda=(1+u)(1+v)/(uv)$ and $\rho=uv/(1+u+v)\in\Z[[u,v]]$. For every
$R(x)=\sum_lr_lx^{2l}\in\Q[[x]]$ and $N\ge1$,
\[
 \CT\bigl(\Lambda^N\rho\,R(uv)\bigr)=\sum_{n<N}\sum_lr_l\binom n{2l}^2 ,
\]
the constant term being taken in $(uv)^{-N}\Q[[u,v]]$. In particular
$H_2(N)=\CT(\Lambda^N\rho\Phi(uv))$.
\end{proposition}
\begin{proof}
We have $\Lambda-1=(1+u+v)/(uv)=\rho^{-1}$, so $\Lambda^N\rho=\sum_{n<N}\Lambda^n+\rho$.
Moreover $\CT(\rho R(uv))=0$, because $\rho\in uv\Z[[u,v]]$. Finally
$\CT(\Lambda^n(uv)^{2l})=[u^{n-2l}](1+u)^n\,[v^{n-2l}](1+v)^n=\binom n{2l}^2$.
\end{proof}

\begin{lemma}[Divergence criterion]\label{lem:divergence}
Let $p$ be a prime and $\Psi=\sum_{(r,s)\neq(0,0)}\psi_{rs}u^rv^s\in\Zp[[u,v]]$ with
$\psi_{rs}\in p^{\min(v_p(r),v_p(s))}\Zp$ for all $(r,s)$. Then $\CT(\Lambda^N\Psi)\in N\Zp$ for every $N\ge1$.
\end{lemma}
\begin{proof}
For $(r,s)\ne(0,0)$ choose $x_{rs},y_{rs}\in\Zp$ with
$x_{rs}r+y_{rs}s=p^{\min(v_p r,v_p s)}$; this is B\'ezout in the discrete valuation ring $\Zp$.
Then $\Psi=\theta_uJ_1+\theta_vJ_2$, where $J_1,J_2\in\Zp[[u,v]]$ have coefficients
$\psi_{rs}x_{rs}p^{-\min}$ and $\psi_{rs}y_{rs}p^{-\min}$. Since $\CT(\theta_u(\cdot))=0$,
\[
\CT(\Lambda^N\theta_uJ_1)=-\CT(\theta_u(\Lambda^N)J_1)=-N\,\CT\bigl(\Lambda^{N-1}\theta_u(\Lambda)J_1\bigr)\in N\Zp,
\]
because $\theta_u\Lambda=-(1+v)/(uv)$ has integer coefficients. The $v$-term is the same.
\end{proof}

\begin{corollary}\label{cor:criterion}
Let $p$ be a prime. Suppose
$[u^rv^s]\,\rho\Phi(uv)\in p^{\min(v_p(r),v_p(s))}\Zp$ for all $r,s\ge0$. Then
$p^{v_p(N)}\mid S(N)$ for all $N\ge1$.
\end{corollary}
\begin{proof}
Combine Propositions~\ref{prop:split} and~\ref{prop:CT} with Lemma~\ref{lem:divergence}.
\end{proof}

We record two more identities used in Sections~\ref{sec:two} and~\ref{sec:CM}:
\begin{align}
 \binom jm^2&=\sum_{i=0}^m\binom mi\binom{m+i}i\binom j{m+i},\label{eq:product}\\
 \sum_{j<n}(20j+17)\binom jt&=(20t+17)\binom n{t+1}+20(t+1)\binom n{t+2}.\label{eq:hockey}
\end{align}
The first counts ordered pairs of $m$-subsets of a $j$-set by the size of their union. The second
follows from the hockey-stick identity and $j\binom jt=t\binom jt+(t+1)\binom j{t+1}$.

\section{The prime 2}\label{sec:two}

The following argument is due to \cite[Theorem 3.1]{K26}.

\begin{theorem}\label{thm:two}
For every $n\ge1$, $S(n)/n$ is an odd $2$-adic integer. In particular $v_2(S(n))=v_2(n)$.
\end{theorem}
\begin{proof}
Using \eqref{eq:product} with $m=2k$ and then \eqref{eq:hockey}, we get, with $t=2k+i$,
\[
\frac{S(n)}n=\sum_k\gamma_k\sum_{i=0}^{2k}\binom{2k}i\binom{2k+i}i
\Bigl(\frac{20t+17}{t+1}\binom{n-1}t+\frac{20(t+1)}{t+2}\binom{n-1}{t+1}\Bigr),
\]
since $\binom n{t+1}/n=\binom{n-1}t/(t+1)$ and $\binom n{t+2}/n=\binom{n-1}{t+1}/(t+2)$. For
$k\ge1$, $v_2(\gamma_k)\ge2k+2$ because $\binom{2k}k$ is even. Also
$t+2\le4k+2<2^{2k+2}$, so every term with $k\ge1$ lies in $2\Z_2$. The term $k=0$ equals
$n^{-1}\sum_{j<n}(20j+17)=10n+7$, which is odd.
\end{proof}

\section{The Hasse invariant and the CM dichotomy}\label{sec:CM}

For an odd prime $p$ and an indeterminate $b$ put
$U_p(b)=\sum_{r<p}\sum_kb^k\binom r{2k}^2\binom{2k}k^2$. Thus $U_p(-4)=\sum_{r<p}a_r$.
Let $P_d$ be the Legendre polynomial.

\begin{proposition}\label{prop:hasse}
Let $p$ be an odd prime. In $\F_p[b,b^{-1}]$,
\[
 d_p(b)\equiv U_p(b)\equiv b^h\sum_{j=0}^dA_j(256b)^{-j}\equiv b^hP_d(t)^2,\qquad t^2=1-\frac1b .
\]
Here $P_d(t)^2$ is a polynomial in $t^2$, so the right side lies in $\F_p[b,b^{-1}]$.
\end{proposition}
\begin{proof}
By \eqref{eq:product} and $\sum_{r<p}\binom r{m+i}=\binom p{m+i+1}$, only
$i=p-1-m$ survives modulo $p$. Hence $\sum_{r<p}\binom r{2k}^2\equiv\binom{2k}{p-1-2k}$, because
$\binom{p-1}{i}\equiv(-1)^i=1$. Therefore
$U_p(b)\equiv\sum_kb^k\binom{2k}k^2\binom{2k}{p-1-2k}$.

Put $k=h-j$ in $U_p$ and $l=h-j$ in \eqref{eq:dN}. The coefficients of $b^{h-j}$ are
$\binom{2h-2j}{h-j}^2\binom{p-1-2j}{2j}$ and $\binom{2h-2j}{h-j}^2\binom{4j}{2j}$ respectively.
For $j\le d$ we have $\binom{p-1-2j}{2j}\equiv\binom{-1-2j}{2j}=\binom{4j}{2j}$, since $(2j)!$ is a
$p$-unit. For $d<j\le h$ both coefficients vanish modulo $p$: $4j>p-1$ forces a carry in
$2j+2j$, and $\binom{p-1-2j}{2j}=0$. So $d_p\equiv U_p$.

Next, for $j\le d$,
$\binom{2h-2j}{h-j}\equiv\binom{-1-2j}{h-j}=(-1)^{h-j}\binom{h+j}{2j}$ and, since $h\equiv-\frac12$,
$\binom{h+j}{2j}\equiv\binom{j-1/2}{2j}=(-1)^j\binom{2j}j16^{-j}$. Hence
$\binom{2h-2j}{h-j}^2\equiv\binom{2j}j^2256^{-j}$, which gives the third expression.

Finally, $A_j256^{-j}=(\frac14)_j(\frac12)_j(\frac34)_j/j!^3$. Also $\{-d,d+1\}\equiv\{\frac14,\frac34\}$,
so $(\frac14)_j(\frac34)_j\equiv(-d)_j(d+1)_j$ for $j\le d$. The classical identity
$P_d(t)^2={}_3F_2(-d,d+1,\frac12;1,1;1-t^2)
=\sum_j\binom dj\binom{d+j}j\binom{2j}j\bigl(\frac{t^2-1}4\bigr)^j$ \cite{Bailey} then gives the last
congruence with $1-t^2=1/b$.
\end{proof}

\begin{proposition}\label{prop:CMdichotomy}
Let $p$ be an odd prime. Then $d_p(-4)\not\equiv0\pmod p$ if and only if $p\equiv1,3,7,9\pmod{20}$.
\end{proposition}
\begin{proof}
For $p=5$, $d_5(-4)=550\equiv0$. Let $p\ne5$. By Proposition~\ref{prop:hasse} at $b=-4$,
$d_p(-4)\equiv(-4)^hP_d(t_0)^2$ with $t_0^2=\frac54$, $t_0\in\overline{\F}_p$. Let $t_0=\sqrt5/2$ and
consider
$E:\ y^2=x\bigl(x^2+2x+\tfrac{1-t_0}2\bigr)$ over $\Q(\sqrt5)$. Its discriminant
$\Delta=8(t_0-1)^2(t_0+1)$ is a unit at every prime above an odd rational prime, because
$(t_0-1)(t_0+1)=\frac14$. Its $j$-invariant is
$64(3t_0+5)^3/((t_0-1)^2(t_0+1))=632000+282880\sqrt5=j(\sqrt{-5})$ \cite{Duke}, so $E$ has complex
multiplication by $\Z[\sqrt{-5}]$, of discriminant $-20$.

The Hasse invariant of the reduction is the coefficient of $x^{p-1}$ in
$(x^3+2x^2+b'x)^h$, with $b'=\frac{1-t_0}2$ \cite[Theorem V.4.1(a)]{Silverman}. Choosing $k$ factors
$b'x$, $h-2k$ factors $2x^2$ and $k$ factors $x^3$, this coefficient is
$2^h\sum_k\frac{(-h/2)_k((1-h)/2)_k}{k!^2}b'^k$. Since $h\equiv-\frac12$, we have
$\{-\frac h2,\frac{1-h}2\}\equiv\{\frac14,\frac34\}\equiv\{-d,d+1\}$, and the sum terminates at
$k=d$. So the coefficient is
$\equiv2^h{}_2F_1(-d,d+1;1;b')=2^hP_d(1-2b')=2^hP_d(t_0)$ (cf.\ \cite[\S5]{K26}; checked in the
Appendix). By Deuring's reduction theorem
\cite{Deuring,EOY}, the reduction at a prime above $p\nmid20$ is ordinary if and only if $p$ splits
in $\Q(\sqrt{-5})$, i.e.\ $\left(\frac{-20}p\right)=1$, i.e.\ $p\equiv1,3,7,9\pmod{20}$.
\end{proof}

\begin{corollary}\label{cor:units}
Let $p\equiv1,3,7,9\pmod{20}$ and $t_0=-1/1024$. Then $d_p(-4)$, $U_p(-4)$ and
$\cF_{<p}(t_0)=\sum_{k<p}A_kt_0^k$ are $p$-adic units.
\end{corollary}
\begin{proof}
$A_k\equiv0\pmod p$ for $d<k<p$, so $\cF_{<p}(t_0)\equiv\sum_{k\le d}A_k(-1024)^{-k}$. This is
$(-4)^{-h}d_p(-4)$ modulo $p$ by Proposition~\ref{prop:hasse}.
\end{proof}

\begin{lemma}[Lucas factorisation]\label{lem:lucas}
For every odd prime $p$ and $s\ge1$, $d_{p^s}(b)\equiv d_p(b)^{(p^s-1)/(p-1)}\pmod p$ in $\Z[b]$.
\end{lemma}
\begin{proof}
Write $l=l_0+pl'$ with $0\le l_0<p$. Lucas' theorem gives
$\binom{2l}l\equiv\binom{2l_0}{l_0}\binom{2l'}{l'}\pmod p$; when $2l_0\ge p$ both sides vanish. In
$\F_p[b]$ we also have $b^l=b^{l_0}(b^{l'})^p$. Hence $G_b(x)\equiv G_{b,<p}(x)\,G_b(x)^p$ in
$\F_p[b][[x]]$, where $G_{b,<p}=\sum_{2l_0\le p-1}b^{l_0}\binom{2l_0}{l_0}^2x^{2l_0}$. Likewise
$E(x)\equiv E_{<p}(x)E(x)^p$ with $E_{<p}=\sum_{2i\le p-1}\binom{2i}ix^i$. So
$\Psi:=G_bE\equiv\Psi_1\Psi^p$ with $\Psi_1=G_{b,<p}E_{<p}$ of degree $<2p-1$. Since
$\Psi^p$ is a series in $x^p$, the coefficient of $x^{p^s-1}$ in $\Psi_1\Psi^p$ comes only from
$x^{p-1}$ in $\Psi_1$. As $[x^{p-1}]\Psi_1=[x^{p-1}]\Psi=d_p(b)$, we get
$d_{p^s}\equiv d_p\,d_{p^{s-1}}^{\,p}$, and the claim follows by induction on $s$.
\end{proof}

\section{The primes 5 and \texorpdfstring{$p\equiv11,13,17,19\pmod{20}$}{inert primes}}\label{sec:inert}

This section reproduces, with complete proofs, the exact digit operator of \cite[\S\S5--6]{K26}.
Fix a prime $p\ge5$ and put $w(n,k)=(-4)^kM(n,k)^2$ for $0\le2k\le n$, and $w(n,k)=0$ otherwise.
Let $\mathcal A_p=\Zp\langle X,Y\rangle$ be the ring of restricted power series. Its coefficients
tend to $0$, so its elements can be evaluated on $\Zp^2$. Define
\[
 I_N(F)=\sum_{n<N}\sum_kw(n,k)F(n,k),\qquad U(N)=I_N(1),\ T(N)=I_N(X),\ V(N)=I_N(Y),
\]
so $S(N)=I_N(20X+17)$, $U(N)=\sum_{n<N}a_n$ and $T(N)=\sum_{n<N}na_n$. Put $W=(p-1)!$,
$E_p(N)=\prod_{1\le i\le N,\,p\nmid i}i$, $H_r=\sum_{i\le r}1/i$ and
$Q=((-4)^{p-1}-1)/p\in\Zp$.

\begin{lemma}\label{lem:unitfactorial}
There are $\mathcal H_r\in1+p\Zp\langle X\rangle$ ($0\le r<p$) with
$E_p(pq+r)=W^qr!\,\mathcal H_r(q)$ for all integers $q\ge0$, and
$\mathcal H_r(X)\equiv1+pH_rX\pmod{p^2}$.
\end{lemma}
\begin{proof}
Let $B(X)=\prod_{i=1}^{p-1}(1+pX/i)$. Then $\log B(X)=\sum_{j\ge1}\ell_jX^j$ with
$\ell_j=(-1)^{j-1}p^j\eta_j/j$, where $\eta_j=\sum_{i<p}i^{-j}$. Let
$S_j(X)=\sum_{a=0}^j\genfrac\{\}{0pt}{}{j}{a}(X)_{\underline{a+1}}/(a+1)$ be the polynomial with
$S_j(q)=\sum_{h<q}h^j$. Its coefficients have valuation $\ge-\lfloor\log_p(j+1)\rfloor$. Put
$\mathcal A(X)=\sum_j\ell_jS_j(X)$. By Wolstenholme, $v_p(\eta_1)\ge2$ and $v_p(\eta_2)\ge1$
for $p\ge5$; for $j\ge3$ we have $j-v_p(j)-\lfloor\log_p(j+1)\rfloor\ge3$, with a bound tending
to $\infty$. Hence $\mathcal A\in p^3\Zp\langle X\rangle$ and
$\mathcal E=\exp\mathcal A\in1+p^3\Zp\langle X\rangle$. At integers $q\ge0$ we get
$\mathcal E(q)=\prod_{h<q}B(h)$. Set $\mathcal H_r(X)=\mathcal E(X)\prod_{i\le r}(1+pX/i)$. Then
$E_p(pq+r)=\prod_{h<q}\prod_{i<p}(ph+i)\cdot\prod_{i\le r}(pq+i)=W^qr!\,\mathcal H_r(q)$.
\end{proof}

\begin{lemma}\label{lem:kernel}
For $0\le r,s<p$ let $\varepsilon=\lceil(2s-r)/p\rceil\in\{0,1,2\}$, $t=r-2s+\varepsilon p\in[0,p)$
and $z=X-2Y$. Put
\[
 K_{rs}=p^{2\varepsilon}(z)_{\underline\varepsilon}^2\frac{(-4)^sW^{2\varepsilon}r!^2}{s!^4t!^2}
 \frac{\mathcal C(Y)\mathcal H_r(X)^2}{\mathcal H_s(Y)^4\mathcal H_t(X-2Y-\varepsilon)^2}\in\mathcal A_p,
 \qquad \mathcal C(Y)=\exp\bigl(Y\log((-4)^{p-1})\bigr).
\]
Then $w(pq+r,pj+s)=w(q,j)K_{rs}(q,j)$ for all $q,j\ge0$, and $K_{rs}(0,0)=w(r,s)$. If
$\varepsilon\ge1$ then $K_{rs}\in p^2\mathcal A_p$. If $\varepsilon=0$ then
$K_{rs}\equiv w(r,s)\bigl[1+p\{2X(H_r-H_t)+Y(Q+4H_t-4H_s)\}\bigr]\pmod{p^2}$.
\end{lemma}
\begin{proof}
Use $N!=p^{\lfloor N/p\rfloor}\lfloor N/p\rfloor!\,E_p(N)$ and Lemma~\ref{lem:unitfactorial} for
$n=pq+r$, $k=pj+s$ and $n-2k=p(q-2j-\varepsilon)+t$. The powers of $W$ cancel down to
$W^\varepsilon$, and $q_0!/(q_0-\varepsilon)!=(q_0)_{\underline\varepsilon}$ with $q_0=q-2j$. Also
$(-4)^{pj+s}=(-4)^{j+s}\mathcal C(j)$. If $0\le q_0<\varepsilon$, both sides vanish, and if $2j>q$ both
sides are $0$. The denominators are units of $\mathcal A_p$, since affine integral substitutions
preserve $1+p\mathcal A_p$. The remaining claims follow from
$\mathcal C\equiv1+pQY$ and $\mathcal H_r\equiv1+pH_rX\pmod{p^2}$.
\end{proof}

\begin{proposition}\label{prop:block}
For $F\in\mathcal A_p$ put $(\mathcal T F)(X,Y)=\sum_{r,s<p}K_{rs}(X,Y)F(pX+r,pY+s)$. Then
$I_{pN}(F)=I_N(\mathcal TF)$ for $N\ge1$. Moreover there are $\alpha_F,\beta_F\in\Zp$ and
$R_F\in\mathcal A_p$ with $\mathcal TF=I_p(F)+p\alpha_FX+p\beta_FY+p^2R_F$, so that
\begin{equation}\label{eq:exactblock}
 I_{pN}(F)=I_p(F)U(N)+p\alpha_FT(N)+p\beta_FV(N)+p^2I_N(R_F).
\end{equation}
\end{proposition}
\begin{proof}
Split $n=pq+r$ and $k=pj+s$ and apply Lemma~\ref{lem:kernel}. For the second statement use
$F(pX+r,pY+s)\equiv F(r,s)+pXF_X(r,s)+pYF_Y(r,s)\pmod{p^2}$, which holds for restricted series by
Taylor's formula with divided derivatives. The constant term is
$\sum_{r,s}w(r,s)F(r,s)=I_p(F)$.
\end{proof}

\begin{theorem}\label{thm:induction}
Suppose $U(p),T(p),V(p)\in p\Zp$. Then for all $e,m\ge1$ and $F\in\mathcal A_p$ we have
$I_{p^em}(F)\in p^{e-1}\Zp$, and $I_{p^em}(F)\in p^e\Zp$ when $I_p(F)\in p\Zp$.
\end{theorem}
\begin{proof}
Induct on $e$. The case $e=1$ follows from \eqref{eq:exactblock} with $N=m$. For $e\ge2$ put
$N=p^{e-1}m$. By the strong assertion at $e-1$, applied to $F=1,X,Y$ (whose $I_p$ are
$U(p),T(p),V(p)\in p\Zp$), we get $U(N),T(N),V(N)\in p^{e-1}\Zp$. By the weak assertion,
$I_N(R_F)\in p^{e-2}\Zp$. Now read off \eqref{eq:exactblock}.
\end{proof}

\begin{theorem}\label{thm:inert}
If $p\equiv11,13,17,19\pmod{20}$, then $p^e\mid S(p^em)$ for all $e,m\ge1$.
\end{theorem}
\begin{proof}
By Propositions~\ref{prop:hasse} and~\ref{prop:CMdichotomy}, $U(p)=U_p(-4)\equiv(-4)^hP_d(t_0)^2=0$.
Next, $V(p)=(\theta_bU_p)(-4)$. Differentiating the polynomial congruence
$U_p(b)\equiv b^hP_d(t)^2$ and using $P_d(t_0)=0$ and $t_0\ne0$ gives $V(p)\equiv0$. Finally,
$F_k(X)=\binom X{2k}^2$ satisfies $F_k(2k-1-X)=F_k(X)$ over $\F_p$. Summing over $X\in\F_p$ gives
$\sum_{r<p}rF_k(r)\equiv(k-\frac12)\sum_rF_k(r)$, hence $T(p)\equiv V(p)-\frac12U(p)\equiv0$.
Apply Theorem~\ref{thm:induction} with $F=20X+17$, using $I_p(F)=S(p)=20T(p)+17U(p)$.
\end{proof}

\begin{theorem}\label{thm:five}
For all $e,m\ge1$, $5^e\mid S(5^em)$.
\end{theorem}
\begin{proof}
Take $p=5$. From Lemma~\ref{lem:kernel},
$\mathcal T(1)=U(5)+5\alpha X+5\beta Y+25R$ with $\alpha=\sum2w(r,s)(H_r-H_t)$ and
$\beta=\sum w(r,s)(Q+4H_t-4H_s)$, the sums over pairs with $\varepsilon=0$. Exact evaluation gives
$U(5)=-155$, $\alpha=1440$ and $\beta=-37270/3$, all divisible by $5$ (see the Appendix). Hence
$\mathcal T(1)=C+25R$ with $C\in5\Z_5$. Now we prove, by induction on $e$:
(a) $I_{5^em}(F)\in5^{e-1}\Z_5$ for all $F$; (b) $U(5^em)\in5^e\Z_5$.
For $e=1$, (a) is integrality and (b) is $U(5m)=CU(m)+25I_m(R)$. For the step, use
\eqref{eq:exactblock} for (a). For (b), use $U(5N)=CU(N)+25I_N(R)$ with $N=5^{e-1}m$.
Finally, $F=20X+17=17+5(4X)$ gives $S(5^em)=17U(5^em)+20T(5^em)\in5^e\Z_5$.
\end{proof}

\section{Dwork crystals}\label{sec:BV}

We recall the results of Beukers and Vlasenko \cite{BV1} that will be used; see also \cite{BV2,BV3}.
Let $p$ be an odd prime and let $R$ be a $p$-adically complete integral domain of characteristic
$0$ with $\bigcap_sp^sR=0$. Let $f\in R[x_1^{\pm1},\dots,x_n^{\pm1}]$ have Newton polytope
$\Delta\subset\mathbb R^n$, and let $\mu=\Delta^\circ$ be its interior. We assume
$\mu\cap\Z^n=\{u_0\}$ is a single point. Let $\Omega_f(\mu)$ be the $R$-module spanned by the
rational functions $(m-1)!\,g/f^m$, where $m\ge1$ and $g\in R[x^{\pm1}]$ is supported in $m\mu$.
Let $\widehat\Omega_f(\mu)$ be its $p$-adic completion. Derivations $\delta$ of $R$ act
coefficientwise and preserve $\widehat\Omega_f(\mu)$ \cite[Prop.~3.4]{BV1}.

If $\nu$ is a vertex of $\Delta$ whose coefficient $f_{[\nu]}$ in $f$ lies in $R^\times$, every
element of $\widehat\Omega_f(\mu)$ has a formal Laurent expansion supported in the cone
$C(\Delta-\nu)$ \cite[\S2]{BV1}. A formal Laurent series $\sum a_kx^k$ is a \emph{formal derivative},
i.e.\ lies in $\sum_i\theta_{x_i}(\cdot)$, if and only if $a_k\in\gcd(k_1,\dots,k_n)R$ for all $k$
\cite[Lemma~2.2]{BV1}. Let $U_f(\mu)$ be the submodule of elements whose expansion is a formal
derivative. It does not depend on the choice of $\nu$. Fix a Frobenius lift $\sigma$ on $R$; all rings below
carry one, e.g.\ $b\mapsto b^p$. Then by \cite[Prop.~4.2]{BV1} the submodule equals
$\{\omega:C_p^s(\omega)\in p^s\widehat\Omega_{f^{\sigma^s}}(\mu)\ \forall s\}$, and the Cartier
operator $C_p$ is given by formulas independent of the vertex \cite[\S3, (4)--(7)]{BV1}. The
\emph{Hasse--Witt invariant} $\beta_p\in R$ is the coefficient of $x^{(p-1)u_0}$ in $f^{p-1}$.

\begin{theorem}[{Beukers--Vlasenko \cite[Theorems 4.3 and 5.7]{BV1}}]\label{thm:BV}
Assume $\beta_p\in R^\times$ and put $\omega_0=x^{u_0}/f$.
\begin{enumerate}[label=(\roman*)]
\item $\widehat\Omega_f(\mu)/U_f(\mu)$ is a free $R$-module of rank one with basis the class of
$\omega_0$. Hence for every derivation $\delta$ of $R$ there is a unique $\eta_\delta\in R$ with
$\delta(\omega_0)-\eta_\delta\,\omega_0\in U_f(\mu)$.
\item Let $\nu$ be a vertex of $\Delta$ with $f_{[\nu]}\in R^\times$, let
$k\in C(\Delta-\nu)\cap\Z^n$, and let $a_m\in R$ be the coefficient of $x^m$ in the expansion of
$\omega_0$ at $\nu$. Then
$\delta(a_{p^sk})\equiv\eta_\delta\,a_{p^sk}\pmod{p^sR}$ for all $s\ge1$.
\end{enumerate}
\end{theorem}

By the vertex-independence of $U_f(\mu)$, the element $\eta_\delta$ is the same for every vertex
used in~(ii). This is the key point of Section~\ref{sec:identification}.

\section{The unit-root crystal attached to \texorpdfstring{$S$}{S}}\label{sec:crystal}

Let $p$ be an odd prime. In the variables $(u,v,z,w)$ put
\[
 f'_b=(1+u+v)\bigl(1-uv(z+b/z)(w+1/w)\bigr).
\]
Let $\mathrm{HW}'(b)\in\Z[b]$ be its Hasse--Witt polynomial (computed below), and let $R_b$ be the
$p$-adic completion of $\Zp[b,b^{-1},\mathrm{HW}'(b)^{-1}]$. This is a $p$-adically complete
domain, since its reduction $\F_p[b,b^{-1},\overline{\mathrm{HW}'}^{-1}]$ is a domain and it has no
$p$-torsion.

\begin{lemma}\label{lem:polytopeB}
\begin{enumerate}[label=(\roman*)]
\item The Newton polytope $\Delta'$ of $f'_b$ has the unique interior lattice point $u_0=(1,1,0,0)$,
and $0$ is a vertex with coefficient $1$.
\item $\mathrm{HW}'(b)\equiv d_p(b)\pmod p$.
\item In the expansion of $\omega_0=uv/f'_b$ at the vertex $0$, the coefficient of
$u^rv^sz^0w^0$ is $[u^rv^s]\,\rho\,G_b(uv)$. In particular the coefficient of
$(uv)^N=x^{Nu_0}$ is $d_N(b)$.
\end{enumerate}
\end{lemma}
\begin{proof}
(i) $\Delta'$ is the Minkowski sum of the triangle $\mathrm{conv}\{0,e_u,e_v\}$ and the pyramid
$\mathrm{conv}\{0,e_u+e_v\pm e_z\pm e_w\}$. So its points are $(\alpha+\mu,\beta+\mu,z,w)$ with
$\alpha,\beta\ge0$, $\alpha+\beta\le1$, $0\le\mu\le1$ and $|z|,|w|\le\mu$. If $(r,s,i,j)$ is an
interior lattice point, moving in the $z$- or $w$-directions forces $|i|,|j|<1$. Moving in the
$u$-, $v$- and $(u+v)$-directions forces $r,s\ge1$ and $r+s<3$. Hence $(r,s,i,j)=(1,1,0,0)$, which
is interior (take $\mu=0.6$, $\alpha=\beta=0.4$). All other monomials of $f'_b$ have positive
$u$- or $v$-degree, so $0$ is a vertex.

(ii) The coefficient of $(uv)^{p-1}z^0w^0$ in $f'^{\,p-1}_b$ is
$\sum_l\binom{p-1}{2l}b^l\binom{2l}l^2\cdot\frac{(p-1)!}{a!^2(p-1-2a)!}$ with $a=p-1-2l$. Modulo $p$,
$\binom{p-1}{2l}\equiv1$, and the multinomial is $\binom{p-1}{2a}\binom{2a}a\equiv\binom{2a}a$ for
$2a\le p-1$. When $2a>p-1$ the multinomial vanishes and so does $\binom{2a}a\bmod p$. This gives
$d_p(b)$.

(iii) $1/f'_b=(1+u+v)^{-1}\sum_j(uv)^j(z+b/z)^j(w+1/w)^j$, and
$\CT_{z,w}(z+b/z)^j(w+1/w)^j=b^{j/2}\binom j{j/2}^2$ for even $j$ (and $0$ otherwise). Also
$[u^iv^i](1+u+v)^{-1}=\binom{2i}i$.
\end{proof}

By Lemma~\ref{lem:polytopeB}, Theorem~\ref{thm:BV} applies to $f'_b$ over $R_b$. Let
$\eta_b\in R_b$ denote $\eta_\delta$ for $\delta=\theta_b$:
\begin{equation}\label{eq:etab}
 \theta_b\Bigl(\frac{uv}{f'_b}\Bigr)\equiv\eta_b\,\frac{uv}{f'_b}\pmod{U_{f'_b}(\mu)} .
\end{equation}

\begin{proposition}\label{prop:onenumber}
Let $p\equiv1,3,7,9\pmod{20}$. Then evaluation $b\mapsto-4$ extends to a continuous ring
homomorphism $R_b\to\Zp$. If $20\,\eta_b(-4)+7=0$, then $p^{v_p(N)}\mid S(N)$ for every $N\ge1$.
\end{proposition}
\begin{proof}
$\mathrm{HW}'(-4)\equiv d_p(-4)$ is a unit by Corollary~\ref{cor:units}, which gives the evaluation map.
Let $\Omega=(20\theta_b+7)(uv/f'_b)\in\widehat\Omega_{f'_b}(\mu)$. By \eqref{eq:etab},
$\Omega-(20\eta_b+7)uv/f'_b\in U_{f'_b}(\mu)$. So every coefficient $a_k(b)\in R_b$ of its
expansion at $0$ lies in $\gcd(k)R_b$, and hence $a_k(-4)\in\gcd(k)\Zp$. Since
$20\eta_b(-4)+7=0$, the coefficients of $\Omega|_{b=-4}$ lie in $\gcd(k)\Zp$. By
Lemma~\ref{lem:polytopeB}(iii), the coefficient of $u^rv^sz^0w^0$ in $\Omega|_{b=-4}$ is
$[u^rv^s]\,\rho\,((20\theta_b+7)G_b)(uv)|_{b=-4}=[u^rv^s]\rho\Phi(uv)$, and
$\gcd(r,s,0,0)=\gcd(r,s)$. Conclude with Corollary~\ref{cor:criterion}.
\end{proof}

\section{Ramanujan's crystal and Sun's theorem}\label{sec:ramanujan}

Let $\cF(T)=\sum_{k\ge0}A_kT^k$, $\cF_{<m}(T)=\sum_{k<m}A_kT^k$, and
$\eta_{\mathcal R}=\theta_T\log\cF(T)\in\Zp[[T]]$. Let $R_T$ be the $p$-adic completion of
$\Zp[T,\cF_{<p}(T)^{-1}]$, with Frobenius lift $\sigma(T)=T^p$. Since $\cF_{<p}(0)=1$, expansion at
$T=0$ gives a map $R_T\to\Zp[[T]]$. It is injective and $p$-saturated: if the expansion of $r$
lies in $p^s\Zp[[T]]$ then $r\in p^sR_T$. Indeed, approximate $r$ modulo $p^s$ by
$g/\cF_{<p}^k$ with $g\in\Zp[T]$; then $g\in p^s\Zp[T]$ because $\cF_{<p}$ is a unit power series.

\begin{proposition}\label{prop:dwork}
(i) $\eta_{\mathcal R}\in R_T$, and $\eta_{\mathcal R}\equiv\theta_T\log\cF_{<p^s}\pmod{p^sR_T}$ for all $s\ge1$.
(ii) If $p\equiv1,3,7,9\pmod{20}$, then evaluation at $t_0=-1/1024$ extends to $R_T\to\Zp$, and
$\eta_{\mathcal R}(t_0)=-\frac3{20}$.
\end{proposition}
\begin{proof}
(i) We have $A_k=\CT(g^k)$ for $g=(1+x_1+x_2+x_3)^4/(x_1x_2x_3)$, whose Newton polytope has $0$ as
its only interior lattice point. By Dwork's congruences in the form of Mellit--Vlasenko \cite{MV}
(see also \cite[\S1]{BV2}),
$\frac{\cF(T)}{\cF(T^p)}\equiv\frac{\cF_{<p^s}(T)}{\cF_{<p^{s-1}}(T^p)}\pmod{p^s\Zp[[T]]}$ for all
$s\ge1$. Lucas congruences give $\cF_{<p^s}\equiv\cF_{<p}^{(p^s-1)/(p-1)}\pmod p$, so the right
sides are units of $R_T$. By saturation they form a Cauchy sequence in $R_T$, whose limit
$\lambda$ has expansion $\cF(T)/\cF(T^p)$. The formal identity
$\eta_{\mathcal R}=\theta_T\log\lambda+p\,\sigma(\eta_{\mathcal R})$ gives
$\eta_{\mathcal R}=\sum_{i\ge0}p^i\sigma^i(\theta_T\log\lambda)\in R_T$.

For the congruence, put $D_s=\theta_T\log\cF_{<p^s}-\eta_{\mathcal R}\in R_T$, with $D_0=-\eta_{\mathcal R}$.
Taking $\theta_T\log$ of Dwork's congruence (a congruence of units modulo $p^s$) gives
$D_s\equiv p\,\sigma(D_{s-1})\pmod{p^s}$. Hence by induction $D_s\in p^sR_T$.

(ii) The evaluation exists because $\cF_{<p}(t_0)\in\Zp^\times$ (Corollary~\ref{cor:units}).
Sun proved \cite[Theorem 1.2, (1.5)]{SunEJC} that for every $n\ge1$,
\[
 4(2n+1)\binom{2n}n\ \Bigm|\ \sum_{k=0}^n(20k+3)\binom{2k}k^2\binom{4k}{2k}(-2^{10})^{n-k}.
\]
Take $n=p^s-1$. Adding $(p^s-1)+(p^s-1)$ in base $p$ produces $s$ carries, so
$v_p\binom{2p^s-2}{p^s-1}=s$ by Kummer's theorem. Dividing by the unit $(-1024)^{p^s-1}$,
$\sum_{k<p^s}(20k+3)A_kt_0^k=\bigl((20\theta_T+3)\cF_{<p^s}\bigr)(t_0)\equiv0\pmod{p^s}$. Since
$\cF_{<p^s}(t_0)$ is a unit, part~(i) gives $20\eta_{\mathcal R}(t_0)+3\equiv0\pmod{p^s}$ for every
$s$.
\end{proof}

\section{Identification of the connection}\label{sec:identification}

\begin{theorem}\label{thm:identification}
Let $p$ be an odd prime. The substitution $T\mapsto1/(256b)$ defines a ring homomorphism
$R_T\to R_b$ compatible with $\theta_b=-\theta_T$, and in $R_b$
\[
 \eta_b(b)=-\frac12-\eta_{\mathcal R}\Bigl(\frac1{256b}\Bigr).
\]
\end{theorem}

The substitution is well defined because
$\cF_{<p}(1/(256b))\equiv b^{-h}d_p(b)\equiv b^{-h}\mathrm{HW}'(b)\pmod p$
(Proposition~\ref{prop:hasse} and Lemma~\ref{lem:polytopeB}), which is a unit of $R_b$. The proof
of the identity occupies the rest of this section.

\subsection{A diagonal model}
In the variables $(x,z,w,y)$ put $g(z,w)=(z+1/z)(w+1/w)$, $L(y)=y+2+1/y$ and
\[
 f_c=\bigl(1-c\,x\,g(z,w)\bigr)\bigl(1-x\,L(y)\bigr).
\]
Let $\mathrm{HW}_c(c)$ be its Hasse--Witt polynomial and $R_c$ the $p$-adic completion of
$\Zp[c,c^{-1},\mathrm{HW}_c(c)^{-1}]$.

\begin{lemma}\label{lem:polytopeC}
\begin{enumerate}[label=(\roman*)]
\item The Newton polytope $\Delta$ of $f_c$ has the unique interior lattice point $e_x=(1,0,0,0)$.
Both $0$ (coefficient $1$) and $\nu_1=(2,-1,-1,-1)$ (coefficient $c$) are vertices.
\item $\mathrm{HW}_c(c)\equiv d_p(c^2)\pmod p$. Its leading coefficient in $c$ is
$\binom{p-1}h^2\equiv1$.
\item The map $\iota:R_b\to R_c$, $b\mapsto c^2$, is an injective ring homomorphism with
$\theta_c\circ\iota=2\,\iota\circ\theta_b$.
\item In the expansion of $\omega_1=x/f_c$ at $0$, the coefficient of $x^Nz^0w^0y^0$ is $d_N(c^2)$.
\end{enumerate}
\end{lemma}
\begin{proof}
(i) $\Delta$ is the Minkowski sum of the pyramid $\{(\mu,z,w):0\le\mu\le1,\ |z|,|w|\le\mu\}$ and
the triangle $\{(\nu,y):0\le\nu\le1,\ |y|\le\nu\}$. An interior lattice point $(r,i,j,l)$ must have
$|i|,|j|,|l|<1$ and $0<r<2$, so it is $e_x$, which is interior (take $\mu=\nu=\frac12$). The
functional $e_x-e_z-e_w-e_y$ attains its maximum on the two summands uniquely, at
$e_x-e_z-e_w$ and $e_x-e_y$. So $\nu_1$ is a vertex, and its coefficient is
$(-c)(-1)=c$.

(ii) The coefficient of $x^{p-1}$ in $f_c^{p-1}$ is
$\sum_{a\ \mathrm{even}}\binom{p-1}a c^a\binom a{a/2}^2\binom{p-1}{e}(-1)^{e}\binom{2e}e$ with
$e=p-1-a$. Reduce modulo $p$ as in Lemma~\ref{lem:polytopeB}(ii). The top term $a=p-1$ gives the
leading coefficient.

(iii) $\mathrm{HW}'(c^2)\equiv\mathrm{HW}_c(c)\pmod p$, so it is a unit in $R_c$. Injectivity
holds modulo $p$ (a map of localized polynomial rings over $\F_p$), hence on completions.

(iv) The expansion at $0$ is $x\sum_j(cxg+xL-cx^2gL)^j=x\,(1-cxg)^{-1}(1-xL)^{-1}$. Here
$\CT_{z,w}(1-cxg)^{-1}=\sum_jc^jx^j\CT(g^j)=G_{c^2}(x)$ and $\CT_y(1-xL)^{-1}=E(x)$.
\end{proof}

Let $\eta_c\in R_c$ be the connection coefficient of $f_c$ for $\delta=\theta_c$
(Theorem~\ref{thm:BV}(i)).

\begin{lemma}\label{lem:etac}
$\eta_c=2\,\iota(\eta_b)$.
\end{lemma}
\begin{proof}
Apply Theorem~\ref{thm:BV}(ii) at the vertex $0$ to $f'_b$ with $k=u_0$, and to $f_c$ with $k=e_x$.
By Lemmas~\ref{lem:polytopeB}(iii) and~\ref{lem:polytopeC}(iv),
$\theta_bd_{p^s}(b)\equiv\eta_bd_{p^s}(b)$ in $R_b$ and
$\theta_c\,d_{p^s}(c^2)\equiv\eta_c\,d_{p^s}(c^2)$ in $R_c$, both modulo $p^s$. Apply $\iota$ to
the first and use $\theta_c\circ\iota=2\iota\circ\theta_b$. Since $d_{p^s}(c^2)$ is a unit
(Lemma~\ref{lem:lucas}), $\eta_c\equiv2\iota(\eta_b)\pmod{p^s}$ for all $s$.
\end{proof}

\subsection{The corner expansion}
\begin{lemma}\label{lem:corner}
Let $k_1=e_x-\nu_1=(-1,1,1,1)\in C(\Delta-\nu_1)$ and $m=p^s$, $h_s=(m-1)/2$. In the expansion of
$\omega_1$ at $\nu_1$, the coefficient of $x^{mk_1}=x^{-m}z^mw^my^m$ is
\[
 \alpha_s(c)=\sum_{l=0}^{h_s}c^{-1-2l}\binom{h_s+l}{2l}^2\binom{2m-1-2l}{2l}.
\]
\end{lemma}
\begin{proof}
Write $(1-cxg)(1-xL)=cxg\cdot xL\cdot(1-(cxg)^{-1})(1-(xL)^{-1})$, with
$g=(1+z^2)(1+w^2)/(zw)$ and $L=(1+y)^2/y$. This gives
\[
 f_c=c\,x^{\nu_1}\,(1+z^2)(1+w^2)(1+y)^2\Bigl(1-\frac{zw}{cx(1+z^2)(1+w^2)}\Bigr)\Bigl(1-\frac{y}{x(1+y)^2}\Bigr).
\]
The monomials $z^2,w^2,y,y^2,zw/x,y/x$ occurring in the last five factors lie in
$(\Delta-\nu_1)\setminus\{0\}$. So these factors are units of the ring of Laurent series supported
in the pointed cone $C(\Delta-\nu_1)$, and the expansion of $\omega_1$ at $\nu_1$ is
\[
 \omega_1=\frac{zwy}{cx}\sum_{\alpha,\beta\ge0}\frac{(zw)^\alpha}{(cx)^\alpha}\frac{y^\beta}{x^\beta}
 (1+z^2)^{-1-\alpha}(1+w^2)^{-1-\alpha}(1+y)^{-2-2\beta}.
\]
The coefficient of $x^{-m}z^mw^my^m$ requires $\alpha+\beta=m-1$. The $z$- and $w$-parts each
give $\binom{-1-\alpha}{(m-1-\alpha)/2}$, and the $y$-part gives $\binom{-2-2\beta}{\alpha}$. With
$\alpha=2l$ (forced, as $m$ is odd), these equal $\pm\binom{h_s+l}{2l}$ and $\binom{2m-1-2l}{2l}$,
and the signs square away.
\end{proof}

\begin{lemma}\label{lem:limit}
For fixed $l\ge0$, as $s\to\infty$,
$\binom{h_s+l}{2l}^2\binom{2p^s-1-2l}{2l}\to\binom{4l}{2l}\binom{2l}l^2256^{-l}=A_l256^{-l}$ in $\Zp$.
\end{lemma}
\begin{proof}
Both binomials are polynomials in $p^s$ with coefficients in $\frac1{(2l)!}\Z[\frac12]$, and
$h_s\to-\frac12$, $2p^s-1-2l\to-1-2l$. The limits are
$\binom{l-1/2}{2l}^2=\binom{2l}l^2/256^l$ (see the proof of Proposition~\ref{prop:hasse}) and
$\binom{-1-2l}{2l}=\binom{4l}{2l}$.
\end{proof}

\subsection{Proof of Theorem~\ref{thm:identification}}
Let $u=1/c$. Let $\cL$ be the ring of two-sided series $\sum_{n\in\Z}e_nu^n$ with $e_n\in\Zp$ and
$e_n\to0$ as $n\to-\infty$. (Products converge, since in each coefficient of a product one
factor's index is $\le n/2$.) Expansion at $u=0$ defines a ring homomorphism
$\exp:\Zp[c,c^{-1},\mathrm{HW}_c^{-1}]\to\cL$, since $\mathrm{HW}_c$ has unit leading coefficient. It
maps $p^s$-multiples to $p^s$-multiples, so it extends continuously to $R_c$, it commutes with
$\theta_c=-\theta_u$, and it is injective by the saturation argument of
Section~\ref{sec:ramanujan}. On the image of $R_T\to R_c$, $T\mapsto1/(256c^2)$, it is given by
$T\mapsto u^2/256$ applied to expansions at $T=0$.

By Theorem~\ref{thm:BV}(ii) at the vertex $\nu_1$ (coefficient $c\in R_c^\times$) and
Lemma~\ref{lem:corner},
\[
 \theta_c\,\alpha_s-\eta_c\,\alpha_s\in p^sR_c,\qquad\text{hence}\qquad
 \theta_c\,\alpha_s-\exp(\eta_c)\,\alpha_s\in p^s\cL\quad(s\ge1).
\]
Here $\alpha_s\in u\Z[u]$. Let $\Pi(u)=u\,\cF(u^2/256)=\sum_lA_l256^{-l}u^{2l+1}$. By
Lemma~\ref{lem:limit}, $\alpha_s\to\Pi$ coefficientwise. Write $\exp(\eta_c)=\sum e_nu^n$. The
coefficient of $u^n$ in $\exp(\eta_c)\alpha_s$ is $\sum_{j\ge1}e_{n-j}[u^j]\alpha_s$, with
$|[u^j]\alpha_s|\le1$ and $e_{n-j}\to0$ as $j\to\infty$. So it converges to
$\sum_je_{n-j}[u^j]\Pi$ as $s\to\infty$. Letting $s\to\infty$ gives
$\theta_c\Pi=\exp(\eta_c)\Pi$ in $\cL$. Since $\Pi\in u(1+u^2\Zp[[u]])$ is invertible in
$u^{-1}\Zp[[u]]\subset\cL$,
\[
 \exp(\eta_c)=\frac{\theta_c\Pi}\Pi=-\theta_u\log\Pi=-1-2\,(\theta_T\log\cF)(u^2/256)
 =\exp\Bigl(-1-2\,\eta_{\mathcal R}\bigl(\tfrac1{256c^2}\bigr)\Bigr).
\]
By injectivity, $\eta_c=-1-2\eta_{\mathcal R}(1/(256c^2))$ in $R_c$. Since
$\iota$ maps $\eta_{\mathcal R}(1/(256b))$ to $\eta_{\mathcal R}(1/(256c^2))$, Lemma~\ref{lem:etac}
turns this into
\[
 \iota\Bigl(\eta_b+\frac12+\eta_{\mathcal R}\bigl(\tfrac1{256b}\bigr)\Bigr)=0 .
\]
Injectivity of $\iota$ finishes the proof.\qed

\section{Proof of the Main Theorem}\label{sec:main}

\begin{proof}[Proof of Theorem~\ref{thm:B}]
Let $p\equiv1,3,7,9\pmod{20}$. Evaluating Theorem~\ref{thm:identification} at $b=-4$, which
corresponds to $T=t_0=-1/1024$, and using Proposition~\ref{prop:dwork}(ii), we get
$\eta_b(-4)=-\frac12+\frac3{20}=-\frac7{20}$. So $20\eta_b(-4)+7=0$, and
Proposition~\ref{prop:onenumber} gives $p^{v_p(N)}\mid S(N)$ for all $N$.
\end{proof}

\begin{proof}[Proof of the Main Theorem]
Let $n\ge1$ and $p\mid n$, and write $n=p^em$ with $p\nmid m$. For $p=2$ use
Theorem~\ref{thm:two}; for $p=5$, Theorem~\ref{thm:five}; for $p\equiv11,13,17,19\pmod{20}$,
Theorem~\ref{thm:inert}; for $p\equiv1,3,7,9\pmod{20}$, Theorem~\ref{thm:B}. Every prime is in
exactly one of these classes, so $p^{v_p(n)}\mid S(n)$ for all $p$, i.e.\ $n\mid S(n)$.
\end{proof}

\begin{remark}
The proof gives more. (a) For every prime $p$ with $p\equiv1,3,7,9\pmod{20}$, $\rho\Phi(uv)$ is a
$p$-adic formal divergence; equivalently,
$p^{\min(v_pr,v_ps)}\mid[u^rv^s]\rho\Phi(uv)$. In particular the diagonal coefficients
$Y_2(n)=[u^nv^n]\rho\Phi(uv)=\sum_l(20l+7)\gamma_l\binom{2i}i$, $i=n-1-2l$, satisfy
$p^{v_p(n)}\mid Y_2(n)$; compare the diagonal criterion of \cite[\S8]{K26}. (b) For every odd prime and every
$b\in\Zp^\times$ with $d_p(b)\not\equiv0$, Theorem~\ref{thm:identification} evaluated at $b$
expresses the unit-root connection of $f'_b$ through Ramanujan's family at $1/(256b)$.
\end{remark}

\section{Discussion}\label{sec:discussion}

\subsection*{Where the new input lies}
Proposition~\ref{prop:split} and Sections~\ref{sec:crystal}, \ref{sec:ramanujan}
and~\ref{sec:identification} are new. Proposition~\ref{prop:CT} is a variant of the
negative-index reduction of \cite{K26}. Sections~\ref{sec:two} and~\ref{sec:inert} reproduce
arguments of \cite{K26}, and the Hasse-invariant computation of Section~\ref{sec:CM} follows
\cite{K26}. We know of no earlier proof of the conjecture. Searches of arXiv and of the literature
on Sun's divisibility conjectures (including \cite{SunEJC,OO,FC}) found the problem listed as open
and no proof. Such a search cannot by itself establish priority. The essential new step is
Theorem~\ref{thm:identification}. It compares two unit-root crystals that are \emph{not} related by
any monomial change of variables: the crystal of $S$, and Ramanujan's crystal at the reciprocal
parameter. Their Hasse invariants agree modulo $p$ (Proposition~\ref{prop:hasse}), but that alone
does not determine a unit root. The comparison comes from the corner expansion of the diagonal
model, whose Hasse--Witt coefficients converge to Ramanujan's period. Conceptually this reflects
the inversion $b\leftrightarrow1/b$ of the hypergeometric equation in Remark~\ref{rem:upside}, and
the Landen-type involution $z\mapsto(1-iz)/(z-i)$, which satisfies
$(z+z^{-1})(\mu(z)+\mu(z)^{-1})=4$ and exchanges the two degenerations of the Legendre family. No
monomial substitution can realise this involution, which is why the comparison has to go through
a limit.

\subsection*{The role of the CM point}
The prime $p$ is split (ordinary) or inert (supersingular) according to the reduction of the CM
curve with $j=j(\sqrt{-5})$. At inert primes the three moments vanish and the digit operator of
Section~\ref{sec:inert} suffices. At split primes the moments are units, and the divisibility is
equivalent to the precise $p$-adic value $\eta_b(-4)=-7/20$. That value is the $p$-adic counterpart
of Ramanujan's formula, supplied by Sun's WZ theorem through Proposition~\ref{prop:dwork}.

\subsection*{Verification}
The Appendix contains an exact-arithmetic script. It checks the conjecture for $n\le2000$ and
every identity, congruence and hypothesis used above on explicit ranges: the decomposition of
Proposition~\ref{prop:split}; the polytope facts; the Hasse--Witt congruences; the agreement of the
two vertex expansions predicted by Theorem~\ref{thm:BV}(ii); Theorem~\ref{thm:identification} at
the level of truncations (82 cases); Sun's theorem; the CM dichotomy for $p<1500$; the moments at
inert primes and at $p=5$; the Legendre--Clausen identity; and the two-variable divisibility of
$\rho\Phi(uv)$. Earlier exact runs verified $n\mid S(n)$ for all $n\le10^6$ \cite{EarlierNote}. No
step of the proof depends on these computations.

\begin{thebibliography}{99}
\bibitem{Bailey} W.~N.~Bailey, \emph{Generalized Hypergeometric Series}, Cambridge Tracts 32,
Cambridge Univ.\ Press, 1935 (Clausen's formula and its terminating cases, \S10).
\bibitem{BV1} F.~Beukers and M.~Vlasenko, Dwork crystals I, \emph{Int.\ Math.\ Res.\ Not.\ IMRN}
2021, no.~12, 8807--8844; arXiv:1903.11155.
\bibitem{BV2} F.~Beukers and M.~Vlasenko, Dwork crystals II, \emph{Int.\ Math.\ Res.\ Not.\ IMRN}
2021, no.~6, 4427--4444; arXiv:1907.10390.
\bibitem{BV3} F.~Beukers and M.~Vlasenko, Dwork crystals III: from excellent Frobenius lifts towards
supercongruences, \emph{Int.\ Math.\ Res.\ Not.\ IMRN} 2023, no.~23, 20433--20483;
arXiv:2105.14841.
\bibitem{Deuring} M.~Deuring, Die Typen der Multiplikatorenringe elliptischer Funktionenk\"orper,
\emph{Abh.\ Math.\ Sem.\ Hansischen Univ.} 14 (1941), 197--272.
\bibitem{Duke} W.~Duke, Modular functions and the uniform distribution of CM points,
\emph{Math.\ Ann.} 334 (2006), 241--252.
\bibitem{EarlierNote} QEEDDD Verifiable-Mathematical-Reasoning-Corpus project, Sun's divisibility
conjecture for $\sum_{k<n}(20k+17)W_k(1,-4)$: squarefree moduli, the second layer, and a reduction
via negative-index continuation, research note, 7 November 2026, \texttt{sun-A179524/paper.tex}.
\bibitem{EOY} N.~Elkies, K.~Ono and T.~Yang, Reduction of CM elliptic curves and modular function
congruences, \emph{Int.\ Math.\ Res.\ Not.} 2005, no.~44, 2695--2707.
\bibitem{FC} Google DeepMind and contributors, \emph{formal-conjectures},
\texttt{FormalConjectures/OEIS/179524.lean},
\url{https://github.com/google-deepmind/formal-conjectures}.
\bibitem{K26} Keio (research note prepared with ChatGPT), Prime-power divisibility in Sun's A179524
conjecture: an integral analytic digit operator and a Catalan diagonal reduction, manuscript,
8 November 2026 (supplied to this project).
\bibitem{MV} A.~Mellit and M.~Vlasenko, Dwork's congruences for the constant terms of powers of a
Laurent polynomial, \emph{Int.\ J.\ Number Theory} 12 (2016), 313--321.
\bibitem{OEIS} OEIS Foundation, The On-Line Encyclopedia of Integer Sequences, sequence A179524
(Z.-W.~Sun, 2010), \url{https://oeis.org/A179524}.
\bibitem{OO} T.~Adamczewski, OEIS Open: How many conjectures can language models turn into
theorems?, arXiv:2608.11941 (2026).
\bibitem{Silverman} J.~H.~Silverman, \emph{The Arithmetic of Elliptic Curves}, 2nd ed., Graduate
Texts in Math.\ 106, Springer, 2009.
\bibitem{SunEJC} Z.-W.~Sun, Products and sums divisible by central binomial coefficients,
\emph{Electron.\ J.\ Combin.} 20(1) (2013), \#P9; arXiv:1004.4623.
\bibitem{Sun2010} Z.-W.~Sun, On Ap\'ery numbers and generalized central trinomial coefficients,
arXiv:1006.2776v11 (2010), Conjecture 5.8.
\end{thebibliography}

\appendix
\section{Verification script and output}\label{app:code}

The following Python~3 script (standard library only) was run as \texttt{python3 verify\_full.py}
(about one minute). Its output is reproduced after the listing.

\begin{lstlisting}
#!/usr/bin/env python3
"""Exact-arithmetic checks for the proof of Sun's A179524 conjecture.

Run:  python3 verify_full.py            (about one to three minutes)
Every check prints PASS or FAIL; the script exits with status 1 on any FAIL.
Nothing here is used as a proof step; the checks confirm each identity,
congruence and hypothesis used in the paper on explicit ranges.
"""
from math import comb, factorial, gcd
from fractions import Fraction as Fr
import sys

FAILED = []


def report(name, ok):
    print(("PASS  " if ok else "FAIL  ") + name)
    if not ok:
        FAILED.append(name)


def primes(n):
    s = bytearray([1]) * (n + 1)
    s[0:2] = b"\x00\x00"
    for i in range(2, int(n ** 0.5) + 1):
        if s[i]:
            s[i * i::i] = bytearray(len(s[i * i::i]))
    return [i for i in range(n + 1) if s[i]]


def gam(l):                      # gamma_l = (-4)^l C(2l,l)^2
    return (-4) ** l * comb(2 * l, l) ** 2


def A(k):                        # Ramanujan coefficients (4k)!/k!^4
    return factorial(4 * k) // factorial(k) ** 4


def a(n):                        # a_n = sum_k (-4)^k C(n,2k)^2 C(2k,k)^2
    return sum(gam(k) * comb(n, 2 * k) ** 2 for k in range(n // 2 + 1))


def vp(x, p):
    if x == 0:
        return 10 ** 9
    v = 0
    while x % p == 0:
        x //= p
        v += 1
    return v


# 1. Sun's conjecture n | S(n), exact, n <= NMAX
NMAX = 2000
S, ok = 0, True
for n in range(1, NMAX + 1):
    S += (20 * (n - 1) + 17) * a(n - 1)
    if S % n:
        ok = False
report("n | S(n) for 1 <= n <= %d" % NMAX, ok)

# 2. Telescoping identity and the decomposition S = 10 H3 + H2 (Prop. 2.4)
ok = True
for m in range(0, 14):
    for N in range(0, 60):
        lhs = sum((2 * n - m + 1) * comb(n, m) ** 2 for n in range(N))
        rhs = N * comb(N, m + 1) * comb(N - 1, m) if N >= 1 else 0
        ok &= (lhs == rhs)
for N in range(1, 300):
    S = sum((20 * n + 17) * a(n) for n in range(N))
    H3 = N * sum(gam(k) * comb(N, 2 * k + 1) * comb(N - 1, 2 * k) for k in range(N))
    H2 = sum((20 * k + 7) * gam(k) * comb(n, 2 * k) ** 2
             for n in range(N) for k in range(n // 2 + 1))
    ok &= (S == 10 * H3 + H2)
report("telescoping identity and S = 10*H3 + H2 (N < 300)", ok)

# 3. Unique interior lattice points of the two Newton polytopes (Lemma 6.4),
#    via the explicit description used in the proof.
def interior_fprime(r, s, i, j):   # exists mu: max(|i|,|j|,(r+s-1)/2) < mu < min(r,s,1)
    lo = max(abs(i), abs(j), Fr(r + s - 1, 2))
    hi = min(r, s, 1)
    return lo < hi


def interior_fc(r, i, j, l):       # exists mu,nu in (0,1): mu+nu=r, mu>max|i|,|j|, nu>|l|
    lo_mu = max(abs(i), abs(j), r - 1)
    hi_mu = min(1, r - abs(l))
    return lo_mu < hi_mu
pts1 = [(r, s, i, j) for r in range(0, 3) for s in range(0, 3)
        for i in range(-2, 3) for j in range(-2, 3) if interior_fprime(r, s, i, j)]
pts2 = [(r, i, j, l) for r in range(0, 3) for i in range(-2, 3)
        for j in range(-2, 3) for l in range(-2, 3) if interior_fc(r, i, j, l)]
report("interior lattice points: f'_b -> %s, f_c -> %s" % (pts1, pts2),
       pts1 == [(1, 1, 0, 0)] and pts2 == [(1, 0, 0, 0)])


def d_coeffs(n):                 # d_n(b) = sum_l b^l C(2l,l)^2 C(2i,i), i = n-1-2l
    return {l: comb(2 * l, l) ** 2 * comb(2 * (n - 1 - 2 * l), n - 1 - 2 * l)
            for l in range((n - 1) // 2 + 1)}


# 4. Hasse-Witt congruences (Lemma 6.4)
ok = True
for p in primes(60)[1:]:
    hw1 = {}
    for l in range((p - 1) // 2 + 1):
        aa = p - 1 - 2 * l
        if 2 * aa <= p - 1:
            mult = factorial(p - 1) // (factorial(aa) ** 2 * factorial(p - 1 - 2 * aa))
            hw1[l] = comb(p - 1, 2 * l) * comb(2 * l, l) ** 2 * mult
    d = d_coeffs(p)
    ok &= all((hw1.get(l, 0) - d.get(l, 0)) % p == 0 for l in set(hw1) | set(d))
    hw2 = {}
    for a2 in range(0, p, 2):
        cc = p - 1 - a2
        hw2[a2] = (comb(p - 1, a2) * comb(a2, a2 // 2) ** 2 * comb(p - 1, cc)
                   * (-1) ** cc * comb(2 * cc, cc))
    ok &= all((hw2.get(e, 0) - (d.get(e // 2, 0) if e % 2 == 0 else 0)) % p == 0
              for e in set(hw2) | {2 * l for l in d})
report("HW(f'_b) == d_p(b) and HW(f_c) == d_p(c^2) mod p, 3 <= p < 60", ok)


def ev(P, x, mod, der=False):
    return sum((e if der else 1) * cf * pow(x, e, mod) for e, cf in P.items()) % mod


def corner(m):                   # corner-vertex coefficients a_{p^s k_1}(c), exponents of c
    h = (m - 1) // 2
    return {-1 - 2 * l: comb(h + l, 2 * l) ** 2 * comb(2 * m - 1 - 2 * l, 2 * l)
            for l in range(h + 1)}


# 5. Two vertex expansions give the same connection (BV Thm 5.7 consistency)
ok, tested = True, 0
for p, s in [(3, 3), (3, 4), (7, 2), (7, 3), (11, 2), (13, 2), (23, 2)]:
    M = p ** s
    mod = M
    cr = corner(M)
    dd = {2 * l: v for l, v in d_coeffs(M).items()}
    for cv in range(2, 14):
        if cv % p == 0:
            continue
        X, D = ev(cr, cv, mod), ev(dd, cv, mod)
        if X % p == 0 or D % p == 0:
            continue
        e1 = ev(cr, cv, mod, True) * pow(X, -1, mod) % mod
        e2 = ev(dd, cv, mod, True) * pow(D, -1, mod) % mod
        ok &= (e1 == e2)
        tested += 1
report("theta log(corner) == theta log(diagonal) mod p^s (%d cases)" % tested, ok)

# 6. Key identity eta_b(b) = -1/2 - eta_R(1/(256 b)) at the level of truncations
ok, tested = True, 0
for p, s in [(3, 4), (7, 3), (11, 2), (13, 2), (23, 2), (29, 2)]:
    M = p ** s
    mod = M
    dd = d_coeffs(M)
    Fc = {k: A(k) for k in range(M)}
    for bv in range(-6, 12):
        if bv % p == 0:
            continue
        D = ev(dd, bv % mod, mod)
        t = pow(256 * bv % mod, -1, mod)
        F = ev(Fc, t, mod)
        if D % p == 0 or F % p == 0:
            continue
        eta = ev(dd, bv % mod, mod, True) * pow(D, -1, mod) % mod
        etaR = ev(Fc, t, mod, True) * pow(F, -1, mod) % mod
        ok &= ((eta + etaR) * 2 + 1) % mod == 0
        tested += 1
report("theta log d_{p^s}(b) + theta log F_{<p^s}(1/(256b)) == -1/2 mod p^s (%d cases)"
       % tested, ok)

# 7. Sun's theorem (EJC 2013, (1.5)) and the instance n = p^s - 1
ok = all(sum((20 * k + 3) * A(k) * (-1024) ** (n - k) for k in range(n + 1))
         % (4 * (2 * n + 1) * comb(2 * n, n)) == 0 for n in range(1, 300))
report("4(2n+1)C(2n,n) | sum_{k<=n}(20k+3)A_k(-1024)^{n-k}, 1 <= n < 300", ok)
ok = True
for q in [3, 9, 27, 81, 7, 49, 23, 29, 41, 43]:
    tot = sum((20 * k + 3) * A(k) * (-1024) ** (q - 1 - k) for k in range(q))
    p = [r for r in primes(q) if q % r == 0][0]
    ok &= (tot % q == 0) and vp(comb(2 * q - 2, q - 1), p) == vp(q, p)
report("p^s | sum_{k<p^s}(20k+3)A_k(-1024)^{p^s-1-k} and v_p C(2p^s-2,p^s-1) = s", ok)

# 8. Ordinary / supersingular dichotomy of the Hasse invariant d_p(-4)
ok = True
for p in primes(1500)[1:]:
    hv = ev(d_coeffs(p), -4 % p, p)
    if p % 20 in (1, 3, 7, 9):
        ok &= hv != 0
    else:
        ok &= hv == 0
report("d_p(-4) != 0 mod p iff p = 1,3,7,9 mod 20 (odd p < 1500)", ok)

# 9. Moments at inert primes and the p = 5 table
def f4(r, s):
    return (-4) ** s * (factorial(r) // (factorial(s) ** 2 * factorial(r - 2 * s))) ** 2
ok = True
for p in primes(400):
    if p % 20 in (11, 13, 17, 19):
        U = sum(f4(r, s) for r in range(p) for s in range(r // 2 + 1))
        T = sum(r * f4(r, s) for r in range(p) for s in range(r // 2 + 1))
        V = sum(s * f4(r, s) for r in range(p) for s in range(r // 2 + 1))
        ok &= (U % p == 0 and T % p == 0 and V % p == 0)
report("U(p), T(p), V(p) == 0 mod p for inert p < 400", ok)
p = 5
H = [Fr(0)] + [sum(Fr(1, i) for i in range(1, r + 1)) for r in range(1, p)]
Q = Fr(4 ** 4 - 1, 5)
al = sum(2 * f4(r, s) * (H[r] - H[r - 2 * s]) for r in range(p) for s in range(r // 2 + 1))
be = sum(f4(r, s) * (Q + 4 * H[r - 2 * s] - 4 * H[s]) for r in range(p) for s in range(r // 2 + 1))
U5 = sum(f4(r, s) for r in range(p) for s in range(r // 2 + 1))


def zero_mod5(x):
    return (x.numerator * pow(x.denominator, -1, 5)) % 5 == 0
report("p = 5: U(5) = %d, alpha = %s, beta = %s, both == 0 mod 5" % (U5, al, be),
       U5 == -155 and zero_mod5(al) and zero_mod5(be))

# 10. Legendre-Clausen identity and the Hasse invariant (Prop. 4.2)
def poly_mul(P, Q2):
    R2 = {}
    for i, x in P.items():
        for j, y in Q2.items():
            R2[i + j] = R2.get(i + j, 0) + x * y
    return R2


def legendre(nn):                # P_n(t) via Bonnet recursion, rational coefficients in t
    P0, P1 = {0: Fr(1)}, {1: Fr(1)}
    if nn == 0:
        return P0
    for k in range(1, nn):
        P2 = {}
        for e, cf in P1.items():
            P2[e + 1] = P2.get(e + 1, 0) + Fr(2 * k + 1, k + 1) * cf
        for e, cf in P0.items():
            P2[e] = P2.get(e, 0) - Fr(k, k + 1) * cf
        P0, P1 = P1, P2
    return P1
ok = True
for dd_ in range(0, 15):
    L = legendre(dd_)
    lhs = poly_mul(L, L)
    rhs = {}
    for j in range(dd_ + 1):
        base = {2: Fr(1, 4), 0: Fr(-1, 4)}
        term = {0: Fr(comb(dd_, j) * comb(dd_ + j, j) * comb(2 * j, j))}
        for _ in range(j):
            term = poly_mul(term, base)
        for e, cf in term.items():
            rhs[e] = rhs.get(e, 0) + cf
    ok &= all(lhs.get(e, 0) == rhs.get(e, 0) for e in set(lhs) | set(rhs))
report("P_d(t)^2 = sum_j C(d,j)C(d+j,j)C(2j,j)((t^2-1)/4)^j, d < 15", ok)
ok = True
for p in primes(60)[3:]:
    h, d = (p - 1) // 2, p // 4
    # coefficient of x^{p-1} in (x(x^2+2x+b'))^h, b'=(1-t)/2, as polynomial in t
    co = {}
    for k in range(d + 1):
        # choose k factors b'x, (h-2k) factors 2x^2, k factors x^3: h!/(k!(h-2k)!k!)
        mult = factorial(h) // (factorial(k) * factorial(h - 2 * k) * factorial(k))
        term = {0: Fr(mult * 2 ** (h - 2 * k))}
        for _ in range(k):
            term = poly_mul(term, {0: Fr(1, 2), 1: Fr(-1, 2)})
        for e, cf in term.items():
            co[e] = co.get(e, 0) + cf
    L = legendre(d)
    for e in set(co) | set(L):
        x = co.get(e, 0) - 2 ** h * L.get(e, 0)
        ok &= (x.numerator * pow(x.denominator, -1, p)) % p == 0
report("Hasse invariant of E_t == 2^h P_d(t) mod p, 7 <= p < 60", ok)

# 11. The two-variable gcd criterion for rho*F(uv) (consequence of the theorem)
ok = True
for r in range(1, 121):
    for s in range(1, 121):
        val = 0
        for l in range((min(r, s) - 1) // 2 + 1):
            aa, bb = r - 1 - 2 * l, s - 1 - 2 * l
            val += (20 * l + 7) * gam(l) * (-1) ** (aa + bb) * comb(aa + bb, aa)
        ok &= (val % gcd(r, s) == 0)
report("gcd(r,s) | [u^r v^s] rho F(uv) for r, s <= 120", ok)

print()
print("ALL CHECKS PASSED" if not FAILED else "FAILED: " + "; ".join(FAILED))
sys.exit(1 if FAILED else 0)
\end{lstlisting}

\noindent Output:
\begin{lstlisting}
PASS  n | S(n) for 1 <= n <= 2000
PASS  telescoping identity and S = 10*H3 + H2 (N < 300)
PASS  interior lattice points: f'_b -> [(1, 1, 0, 0)], f_c -> [(1, 0, 0, 0)]
PASS  HW(f'_b) == d_p(b) and HW(f_c) == d_p(c^2) mod p, 3 <= p < 60
PASS  theta log(corner) == theta log(diagonal) mod p^s (62 cases)
PASS  theta log d_{p^s}(b) + theta log F_{<p^s}(1/(256b)) == -1/2 mod p^s (82 cases)
PASS  4(2n+1)C(2n,n) | sum_{k<=n}(20k+3)A_k(-1024)^{n-k}, 1 <= n < 300
PASS  p^s | sum_{k<p^s}(20k+3)A_k(-1024)^{p^s-1-k} and v_p C(2p^s-2,p^s-1) = s
PASS  d_p(-4) != 0 mod p iff p = 1,3,7,9 mod 20 (odd p < 1500)
PASS  U(p), T(p), V(p) == 0 mod p for inert p < 400
PASS  p = 5: U(5) = -155, alpha = 1440, beta = -37270/3, both == 0 mod 5
PASS  P_d(t)^2 = sum_j C(d,j)C(d+j,j)C(2j,j)((t^2-1)/4)^j, d < 15
PASS  Hasse invariant of E_t == 2^h P_d(t) mod p, 7 <= p < 60
PASS  gcd(r,s) | [u^r v^s] rho F(uv) for r, s <= 120

ALL CHECKS PASSED
\end{lstlisting}

\end{document}
